% =====================================================================
%  Methods chapter
%  Sampling bias in Norwegian biodiversity occurrence data
%
%  Self-contained \chapter to be \input{} into the main thesis.
%  Assumes the thesis preamble already loads, at minimum:
%     \usepackage{booktabs}   % nice tables
%     \usepackage{amsmath}    % math environments
%     \usepackage{hyperref}   % optional
%
%  Code is not reproduced here. Each place that refers to a specific
%  script names it inline as "github code <filename>"; those scripts
%  are published in the accompanying repository.
% =====================================================================

\chapter{Methods}
\label{ch:methods}

\section{Overview: one connected pipeline}
\label{sec:overview}

This chapter describes how the spatial structure of \emph{sampling bias}
in Norwegian species-occurrence data was quantified and, for vascular
plants, validated against an independent structured survey. The analysis
does not refit the occurrence model. It takes as input the posterior
\emph{sampling-intensity} surfaces (the ``bias field'') produced by an
upstream Bayesian spatial model, and interrogates their regional
structure with a sequence of steps that hand off from one to the next.

The logic of the pipeline is deliberately linear. First, the bias field
is the raw object of study, so it is rendered as full-country maps for
visual inspection (Section~\ref{sec:fullmaps}). To turn a visual
impression of spatial structure into something testable, Norway is then
partitioned into a fixed set of small study regions
(Section~\ref{sec:regions}). The field's values are pulled out of those
regions into a single tidy table, giving one observation per raster cell
per group per region (Section~\ref{sec:extract}). That table is then
interrogated by four complementary tests, each asking a different
question of the same extracted values: a \emph{stratified block sampling
test} that localises which region drives the spatial spread
(Section~\ref{sec:block}); a \emph{permutation test} that asks whether
the taxonomic groups are spatially exchangeable
(Section~\ref{sec:permutation}); a \emph{two-way ANOVA} that decomposes
the variance into region, group and interaction components
(Section~\ref{sec:anova}); and a \emph{structured-survey validation} that
checks the bias-corrected predictions against independent field data
(Section~\ref{sec:ano}). The first three tests describe the internal
structure of the modelled bias; the fourth asks whether the correction
built on that bias is trustworthy in the real world.

All steps are implemented in \textsf{R} using \textsf{terra} and
\textsf{sf} for the geospatial operations. The mapping, extraction and
block-sampling steps are organised as a small numbered pipeline driven by
a single configuration file and a shared helper file (github code
\texttt{Workflow/00\_config.R}, github code \texttt{Workflow/01\_setup.R}),
so that switching to a new dataset requires editing only the
configuration; the modules can be run end-to-end through driver scripts
in 6-region and 12-region variants (github code \texttt{Workflow/run\_all.R},
github code \texttt{Workflow/run\_all\_12.R}).

\paragraph{Datasets and the bias field.}
Throughout, the analysis works with several datasets --- \emph{birds},
\emph{fungi}, \emph{vascular plants}, and a re-run bird dataset
(\emph{new birds}). Each dataset is a collection of model \emph{groups}
(taxonomic sub-models), supplied as one posterior raster stack per group
on disk under the convention
\path{<root>/<subfolder>/<prefix><N>/Bias/Bias.rds}. The number of groups
is detected from disk rather than hard-coded, which in the runs reported
here gives on the order of twenty groups per dataset. Each stack is read,
unpacked from its packed \textsf{terra} form, reduced to its
posterior-\emph{mean} layer (the ``bias mean''), and cropped and masked to
the Norwegian mainland so that only on-land cells contribute (github code
\texttt{Workflow/01\_setup.R}). The bias mean is interpreted as a relative
sampling-intensity surface: higher values mark areas the upstream model
infers to be more intensively surveyed.

\paragraph{Coordinate reference and resolution.}
All modelled rasters carry the projection EPSG:25833 (ETRS89 / UTM
zone~33N), but their coordinates are stored in \emph{kilometres} rather
than metres, at an effective resolution of $\approx 1$\,km. This is stated
explicitly because every spatial operation below relies on it: region
squares are defined in this kilometre space, and the point data introduced
for the validation (Section~\ref{sec:ano}) are first reprojected to
EPSG:25833 in metres and then divided by $1000$ before extraction, so that
points and rasters share one coordinate space.

\section{Visualization of the bias-field full maps}
\label{sec:fullmaps}

The first task is simply to see the object under study. For each dataset,
every group's bias-mean surface is plotted as a full-country raster over
the Norwegian outline (github code \texttt{Workflow/02\_full\_maps.R}). The
key design choice is that \emph{all} groups within a dataset are drawn on a
single shared colour scale: a common \texttt{zlim} and a common set of
class breaks are computed across every group first, then reused for each
map (github code \texttt{Workflow/01\_setup.R}). Without this, each map
would be auto-scaled to its own range and the maps would not be visually
comparable; with it, a given colour means the same sampling intensity in
every panel, so differences between groups and between parts of the
country are read directly off the colour.

The full maps are written as a single multi-page PDF (one page per group),
and equivalent one-PNG-per-group versions are produced for dropping
individual maps into figures and slides (github code
\texttt{Workflow/06\_png\_maps.R}; a dataset-general, auto-discovering
variant is github code \texttt{Workflow/06\_png\_maps\_multi.R}). A
companion layout draws, for each group, the full map beside a grid of
zoomed panels --- one per study region --- on the same shared scale, which
previews how the regions of Section~\ref{sec:regions} sit on the field
(github code \texttt{Workflow/03\_region\_maps.R}). As context for what
this modelled intensity is correcting, raw georeferenced occurrence records
for a few exemplar species are also mapped straight from a citable GBIF
download, to show the uneven field effort that motivates the bias modelling
in the first place (github code \texttt{gbif\_norway\_maps.R}).

The purpose of this step is diagnostic and motivating rather than
inferential. The maps make the central pattern visible --- a strong
south--north gradient in inferred sampling intensity, dense in the
well-surveyed south-east and sparse in the far north --- and that visible
pattern is exactly the thing the remaining sections make quantitative.
A map, however, cannot by itself say whether that gradient is statistically
real, whether it is shared across taxonomic groups, or which areas drive
it. To answer those questions the country must be partitioned into
comparable spatial units.

\section{Generating the study regions}
\label{sec:regions}

To compare the bias field across space, Norway is summarised over a set of
fixed $10\times10$\,km square blocks. Each region is defined by a centre
coordinate in the rasters' UTM-33 kilometre space and a side length of
$10$\,km (half-side $5$\,km); square polygons are constructed directly from
those centres in the raster CRS (github code \texttt{Workflow/01\_setup.R}).
Using fixed-size squares rather than, say, administrative units keeps the
spatial ``sample size'' (the block of cells) comparable from region to
region, which the later variance comparisons require.

Two nested region sets are used:

\begin{itemize}
  \item a \textbf{6-region} set --- Setesdal, Oslo, Valdres, Trondheim,
        Tromsø and Lakselv (github code \texttt{Workflow/00\_config.R});
        and
  \item a \textbf{12-region} set --- the original six plus Bergen,
        Kristiansand, Skorovatn, Bodø, Svolvær and Kirkenes (github code
        \texttt{00\_config\_12\_updated.R}; the unadjusted variant is
        github code \texttt{Workflow/00\_config\_12.R}).
\end{itemize}

\noindent
The two sets are intentionally nested: the six regions are exactly the
first six of the twelve, so a single extraction can serve both, and any
result computed on six regions can be compared against the same result on
twelve without re-reading the rasters (used directly in
Section~\ref{sec:anova}). The 6-region set is the original coarse coverage;
the 12-region set roughly doubles spatial coverage and resolution.

\paragraph{Distribution across Norway.}
The centres are chosen to span the country's full survey-effort gradient
rather than to tile it uniformly. They run from the data-rich south-east
(Oslo, Kristiansand) and south (Setesdal, Valdres, Bergen) through the
central transition (Trondheim, Skorovatn), up the northern coast (Bodø,
Svolvær, Tromsø) to the data-sparse far north and Finnmark (Lakselv,
Kirkenes). Spreading the regions along this south--north axis is what lets
the later tests probe whether the bias field's structure tracks the
gradient seen in the maps.

\paragraph{Tuning the northern squares.}
Two of the northernmost squares originally straddled the coastline, so part
of each square fell on ocean cells with no bias value, leaving the regions
with fewer usable cells than the others and unbalancing the design. The
centres of Svolvær and Kirkenes were therefore nudged by a few kilometres
(Svolvær $3$\,km west, Kirkenes $1$\,km south) until each square overlapped
a full, dense block of $121$ non-missing bias cells ($11\times11$ at
$\approx 1$\,km resolution) for every group. The replacement centres were
found by a small search over candidate centres on the raster grid and then
verified to give $121$ valid cells for all groups (github code
\texttt{fix\_region\_centers.R}, github code \texttt{verify\_centers.R}; the
adopted values are recorded in github code \texttt{00\_config\_12\_updated.R}).
This balancing is what makes the variance and ANOVA comparisons fair across
regions.

As a check on interpretation, a separate diagnostic confirms how the
underlying spatial mesh of the upstream model relates to these small
squares: for each $10$\,km region it counts the mesh nodes falling inside
the square and the distance to the nearest node, tabulates them, and draws
a zoom panel per region (github code \texttt{meshRegionPanels.R}; the mesh
itself is rebuilt in github code \texttt{meshNodes.R}). This makes explicit
that regional values are read off a field whose nodes are sparse at the
$10$\,km scale --- relevant context for how much weight to place on any
single region.

The regions define \emph{where} the field is measured; the next step turns
those polygons into actual numbers.

\section{Extracting the values}
\label{sec:extract}

All four tests share one data layout: \emph{one observation per raster
cell, per group, within a region}. For each group, the bias-mean surface
(already cropped and masked to Norway) is intersected with the region
polygons, and every raster cell that \emph{touches} a region square is
extracted together with its cell index, coordinates and bias value (github
code \texttt{Workflow/04\_extract\_values.R}). ``Touching'' cells are used
rather than only cells whose centroid falls inside, so that a square of
fixed geographic size yields a consistent block of cells regardless of how
it happens to align with the grid.

A subtlety matters for the comparisons that follow: every group must be
read on \emph{exactly the same} cells, or apparent differences between
groups could be artefacts of slightly different cell sets. The test scripts
therefore first derive a set of \emph{canonical} cell indices per region
from the reference (Norway-masked) grid, and then align each group's values
to those indices, dropping any cell with a missing value (this
canonical-cell construction lives inside github code
\texttt{permutation\_all\_datasets.R} and github code
\texttt{two\_way\_anova\_all\_datasets.R}). The result is a balanced
long-format table (group $\times$ region $\times$ cell $\rightarrow$ bias
mean), written to a multi-sheet Excel workbook for re-use and cell-count
checking (github code \texttt{Workflow/04\_extract\_values.R}). With the
balanced $121$-cell blocks and roughly twenty groups across twelve regions,
each dataset contributes on the order of $25\,000$--$30\,000$ cell
observations.

A single scalar summary of this table --- the \emph{total variance} of all
pooled, non-missing bias values across every region and group --- is also
computed, and recomputed after the northern centres were re-tuned (github
code \texttt{total\_variance.R}, github code \texttt{new\_total\_variance.R}).
This pooled total variance is the observed statistic the permutation test
(Section~\ref{sec:permutation}) reproduces under its null. With the extracted
table in hand, the four tests can now interrogate it from different angles.

\section{Stratified block sampling test}
\label{sec:block}

The first test asks a localisation question: \emph{which region carries the
most leverage over the spatial spread of the bias field?} The statistic of
interest is the \emph{variance among regional variances}. For a given
group, the bias variance is computed separately within each region, and the
variance of those per-region variances measures how unevenly bias is
structured across space --- a large value means some regions are far more
internally variable than others (github code
\texttt{Workflow/05\_randomization.R}; the script names this a
``randomization'' test, but it is a stratified block sampling test, and is
referred to as such here).

The test proceeds, for each \emph{target} group in turn, as a series of
single-region block swaps:

\begin{enumerate}
  \item \textbf{Original.} Using only the target group's own cells, compute
        the per-region variances and record the variance among them.
  \item \textbf{Single-region swaps.} For each region $r$ and each other
        (\emph{donor}) group $d$, replace region $r$'s block of cells with
        the same region's block taken from donor $d$, while keeping the
        other eleven regions from the target group, and recompute the
        variance among regional variances for this ``mixed'' configuration.
  \item \textbf{Change scores.} For every swap, record the signed and
        absolute change in the statistic relative to the target group's
        original value.
\end{enumerate}

\noindent
Because each swap replaces an entire spatial block (a whole region's cells)
rather than individual cells, the procedure is stratified by region and
operates on coherent blocks --- hence stratified block sampling. Aggregating
the absolute changes \emph{by swapped region} (mean, median, standard
deviation and maximum across all target/donor combinations) gives a
\emph{region-sensitivity} ranking: regions whose replacement typically moves
the statistic the most are the most influential. The outputs are a
per-target summary table, the full table of regional variances, the ten
largest changes per target group, and the region-sensitivity ranking, with
diagnostic plots (a boxplot of absolute change per region, a histogram of
signed changes centred on zero, and a dot chart of mean/median/maximum
change per region). The expectation is that sparsely sampled northern
regions carry the most leverage, since their bias values are the least
constrained.

This test pinpoints \emph{where} structure lives, but it treats one region
at a time and says nothing about whether the taxonomic groups as a whole are
interchangeable. That global question is the permutation test's job.

\section{Permutation test}
\label{sec:permutation}

The permutation test asks whether the taxonomic groups within a dataset are
spatially \emph{exchangeable}: would the pooled bias field look the same if,
within each region, the contributing groups were resampled? If groups carry
idiosyncratic spatial signal, reshuffling which groups contribute should
change the pooled variance; if they do not, the observed total variance
should sit in the middle of the reshuffled distribution. The procedure is
run identically across datasets (github code
\texttt{permutation\_all\_datasets.R}; a single-dataset variant is github
code \texttt{permutation\_histogram.R}):

\begin{enumerate}
  \item For each region, assemble a cells~$\times$~groups matrix
        ($121$ cells $\times$ the detected number of groups $G$) of
        bias-mean values.
  \item \textbf{Observed statistic.} Use each group exactly once, pool all
        regions, and compute the total variance of the pooled values. This
        is the \emph{baseline} --- the same quantity computed in
        Section~\ref{sec:extract}.
  \item \textbf{Null distribution.} Repeat $B=1000$ times: within each
        region independently, draw $G$ groups \emph{with replacement}, take
        that region's $121$ cell values from each drawn group, pool all
        regions, and compute the total variance.
  \item \textbf{Position of the observed value.} Summarise where the
        baseline falls among the $1000$ permuted variances by the proportion
        at least as large,
        $p_{\ge}=\tfrac{1}{B}\sum_b \mathbf{1}\{v_b \ge v_{\mathrm{obs}}\}$
        (and the complementary $p_{\le}$), and plot the null distribution
        with the baseline marked.
\end{enumerate}

\noindent
Resampling \emph{with replacement within region} is the crucial design
choice: it preserves each region's spatial footprint and sample size while
breaking the identity of which group supplied each cell, so the null
isolates group exchangeability from regional structure. The random-number
generator is seeded (\texttt{set.seed(1234)}) so the permutation
distribution is exactly reproducible. A baseline sitting near the centre of
the null ($p_{\ge}\approx 0.5$) is read as evidence that groups are
approximately exchangeable, i.e.\ that the regional structure of the bias
field is largely \emph{shared} across taxonomic groups rather than
group-specific.

Where the block sampling test localises leverage to a region, the
permutation test makes a global statement about groups; together they
separate ``which region matters'' from ``do groups matter at all''. Neither,
however, quantifies \emph{how much} of the field's variation is attributable
to region versus group. That apportionment is what the ANOVA provides.

\section{ANOVA test}
\label{sec:anova}

The two-way ANOVA decomposes the bias field into its two design factors. Per
dataset, the bias mean of each cell is the response, with crossed fixed
factors

\begin{equation}
  y_{ijk} \;=\; \mu \;+\; \alpha_i \;+\; \beta_j \;+\;
  (\alpha\beta)_{ij} \;+\; \varepsilon_{ijk},
  \label{eq:anova}
\end{equation}

\noindent
where $\alpha_i$ is the effect of taxonomic \emph{group}
($i=1,\dots,G$), $\beta_j$ the effect of \emph{region}
($j=1,\dots,12$), $(\alpha\beta)_{ij}$ their interaction, and
$\varepsilon_{ijk}$ the cell-level residual (github code
\texttt{two\_way\_anova\_all\_datasets.R}; a fungi-only version is github
code \texttt{two\_way\_anova\_fungi.R}). Each raster cell within a region,
for a given group, is one observation --- the same long-format layout used
by the other tests (Section~\ref{sec:extract}).

Because the canonical-cell construction makes the design balanced (an equal
number of cells per group$\times$region), the sequential (Type~I) sums of
squares equal the Type~II/III sums of squares and term order is immaterial;
a Type~III check using sum-to-zero contrasts is nonetheless reported where
the \textsf{car} package is available.

\paragraph{Effect sizes, not $p$-values.}
With tens of thousands of cells every term is overwhelmingly
``significant'', so the conclusion is based on \emph{effect size}. For each
term the proportion of variance explained is reported as both $\eta^2$ and
partial $\eta^2$,

\begin{equation}
  \eta^2_{\text{term}} \;=\;
  \frac{\mathrm{SS}_{\text{term}}}{\mathrm{SS}_{\text{total}}},
  \qquad
  \eta^2_{p,\text{term}} \;=\;
  \frac{\mathrm{SS}_{\text{term}}}
       {\mathrm{SS}_{\text{term}}+\mathrm{SS}_{\text{residual}}},
  \label{eq:eta2}
\end{equation}

\noindent
and the dominant systematic source --- region, group or interaction --- is
the term with the largest $\eta^2$. Model fit is checked with the standard
residual diagnostics and a group$\times$region interaction plot of cell
means. The response is modelled on its natural scale; because the intensity
surface is right-skewed, the code also allows the ANOVA to be re-run on
$\log(\text{bias})$ via a switch, and the $\eta^2$ ranking (which factor
dominates) is robust to that choice. The per-dataset decompositions are
collated into a stacked bar chart --- region / group / interaction /
residual share, one bar per dataset --- so the governing factor is visible
at a glance (github code \texttt{anova\_effectsize\_barchart.R}).

\paragraph{Robustness to the number of regions.}
To confirm the decomposition does not depend on how finely Norway is
partitioned, the same two-way ANOVA is refitted on both the 6-region and the
12-region set for each dataset. Because the six regions are nested in the
twelve, each raster is extracted once and both fits come from the same cell
table; the two decompositions are compared in a grouped bar chart so any
shift in the region/group/interaction/residual split is apparent (github
code \texttt{two\_way\_anova\_region\_comparison.R}).

The permutation test, block sampling test and ANOVA together describe the
\emph{internal} structure of the modelled bias --- how it is distributed
across regions and groups, and which regions move it. None of them, however,
can show that the bias-\emph{corrected} predictions built on this field are
actually right. That requires confronting the predictions with data the
model never saw.

\section{Structured-survey test}
\label{sec:ano}

The final component validates the bias-corrected vascular-plant predictions
against an independent, structured survey: the Norwegian nature-monitoring
programme ANO (\emph{Arealrepresentativ naturovervåking}). The validation is
restricted to vascular plants for a substantive reason. ANO records every
vascular plant rooted in each $1\,\text{m}^2$ plot, so a modelled species
\emph{not} listed for a plot is a \emph{true} absence. Birds (mobile,
imperfectly detected) and fungi (no exhaustive Norwegian structured survey)
lack this complete-list property and so cannot supply genuine absences;
their exclusion is itself evidence of the data gaps the thesis documents
(github code \texttt{ano\_validation/00\_config\_ano.R}).

\paragraph{Acquiring and preparing the survey.}
ANO vascular-plant records are obtained from GBIF through a reproducible,
citable download (filtered to the ANO dataset, mainland Norway, the
vascular-plant clade, georeferenced records without geospatial issues, and
coordinate uncertainty $\le 100$\,m), mirroring the upstream project's
filter (github code \texttt{ano\_validation/01\_download\_ano.R}). The
records are cleaned to species level, reprojected to EPSG:25833 in metres,
and clipped to the mainland. The plot unit is the GBIF \texttt{parentEventID}
(year:site:plot), and a plot~$\times$~species presence/absence matrix is
built; because all records carry \texttt{occurrenceStatus = PRESENT},
absences are inferred under the complete-list assumption --- valid precisely
because ANO plots are exhaustive for vascular plants (github code
\texttt{ano\_validation/02\_prepare\_ano.R}). To obtain a genuinely
\emph{independent} test set, the most recent ANO season (2024) is held out;
since ANO rotates roughly one fifth of its sites each year, the held-out
plots are at new locations, under the stated assumption that the 2024 season
had not yet entered the model's training snapshot.

\paragraph{Scoring: discrimination, regional skill, and a bias diagnostic.}
Each modelled species' prediction raster (the posterior-mean occupancy
layer) is read once and extracted at all hold-out plot coordinates --- after
the metres-to-kilometres conversion of Section~\ref{sec:overview} --- and
reused for three summaries (github code
\texttt{ano\_validation/03\_validate\_vascular.R}):

\begin{enumerate}
  \item \textbf{Per-species discrimination.} For each species present in
        both the model and the hold-out, performance is the area under the
        ROC curve (AUC), a threshold-free measure of how well predicted
        occupancy separates present from absent plots. Species with fewer
        than three presences or three absences are not scored; overall
        performance is summarised by mean and median AUC and the fraction of
        species with $\mathrm{AUC}>0.7$.
  \item \textbf{Regional skill.} Each hold-out plot is assigned to the
        nearest of the twelve region centres (the same regions as the rest
        of the pipeline), and a mean AUC is computed per region (regions with
        fewer than $30$ plots are not reported), to show \emph{where} in
        Norway the corrected predictions discriminate well.
  \item \textbf{Bias diagnostic.} A successful correction should not fail
        systematically where the model itself flags high sampling bias. The
        group-level bias field is averaged across segments and sampled at
        each plot; per-plot predictive error is the mean Brier score across
        species ($\overline{(\hat p - y)^2}$, lower is better); and the
        Spearman rank correlation between per-plot bias magnitude and per-plot
        Brier score is reported. A correlation near zero indicates the
        correction holds across the bias gradient, whereas a strong positive
        correlation would signal residual failure in high-bias areas.
\end{enumerate}

\noindent
This closes the loop opened by the maps: the same regional gridding and
bias field that the internal tests dissect are here used to ask whether the
correction actually pays off against independent ground truth, and
specifically whether it pays off \emph{evenly} across the south--north
gradient rather than only where data were already abundant. The
per-species AUCs, per-region AUCs and per-plot bias/error values are written
out for the figures and discussion that follow. Because the saved
predictions are at $\approx 1$\,km resolution, the validation is necessarily
at that resolution; this, the species-coverage limit, the low-support
regions, and the 2024-hold-out assumption are the principal caveats carried
into the Results and Discussion.
